\documentclass[11pt]{article}
\usepackage[hmargin=1.00in,top=.90in,bottom=.88in]{geometry}
\usepackage{amsmath,amssymb,amsthm,mathtools,booktabs,array}
\usepackage[colorlinks=true,linkcolor=blue!60!black,urlcolor=blue!60!black,citecolor=blue!60!black]{hyperref}
\usepackage{xcolor}
\hypersetup{pdftitle={Which parts of an interactad are corners and bispans?},pdfsubject={A relative corners and decorated finite-bispan comparison}}
\newtheorem{theorem}{Theorem}[section]
\newtheorem{lemma}[theorem]{Lemma}
\newtheorem{proposition}[theorem]{Proposition}
\theoremstyle{definition}
\newtheorem{example}[theorem]{Test}
\newcommand{\Set}{\mathsf{Set}}
\newcommand{\Fin}{\mathsf{Fin}}
\newcommand{\Cat}{\mathsf{Cat}}
\newcommand{\V}{\mathcal V}
\newcommand{\J}{\mathcal J}
\newcommand{\Lcat}{\mathcal L}
\newcommand{\Kcat}{\mathcal K}
\newcommand{\Acat}{\mathcal A}
\newcommand{\Rcat}{\mathcal R}
\newcommand{\Bisp}{\mathsf{Bispan}}
\newcommand{\id}{\mathrm{id}}
\newcommand{\bx}{\mathop{\boxtimes}}
\newcommand{\St}{\mathbin{\star}}
\newcommand{\op}{\mathrm{op}}
\newcommand{\Sq}{\mathsf{Sq}}
\newcommand{\hop}{\mathrm{hop}}
\newcommand{\iso}{\xrightarrow{\sim}}
\newcommand{\BC}{\mathrm{BC}}
\newcommand{\ssum}{\mathop{\textstyle\bigoplus}}
\title{Which parts of an interactad are corners and bispans?}
\author{Contextual Composition project technical note}
\date{20 November 2026}
\begin{document}
\maketitle
\noindent This project note was developed with substantial AI assistance.
It is a written comparison with the existing proof, not a Lean certificate.

\noindent\textbf{Verdict.}
The two published constructions do not compose as stated. There is, however,
an explicit factorization of the \emph{presented} crossing. Its geometry is
the chosen codomain-cartesian square of a contentful arrow along a display,
followed by covariant transport along an arbitrary map. Its coefficient
functor is finite polynomial evaluation, but its native arrows live in a
Grothendieck category built from the additive transfers, pulled back along
the augmentation slice. An indexed distributivity isomorphism lifts the
occurrence calculation to that category. The two constructions meet along
the \emph{same addressed carrier}. On the native generators, pure balances
and four structural comparisons their assembly agrees with the given law.
The ten equations then compare the same pastings, not just the same
operations on objects. This is a concrete recovery with a relative
corners lemma and a decorated-family lemma, not a universal
identification of two ambient theories. Marked recognition is separate.
Two shortcuts are false: the augmented fibres need not form a bispan model
with the original products, and restricting the vertical category to
surjections does not give the required relative corners.

\paragraph{Input, size, and meaning of recovery.}
Work in nested universes with ordinary sets. Write
$TB=\coprod_s B^{E_s}$ for the finitary cartesian polynomial monad,
and use the hypotheses for independent interaction in \cite{O1}.
Thus $p:\Lcat\to\Set$ is a global opfibration, cartesian over $r\in\J$
\emph{against arbitrary test arrows}; $d:\Lcat\to\Set$ has natural
$\nu:d\Rightarrow Tp$. Write $c_{r,Q}:r^*Q\to Q$,
$\kappa_{g,Q}:Q\to g_!Q$. Display means
\begin{equation}
 D_{r,Q}=(d c_{r,Q},\nu_{r^*Q}):d(r^*Q)\iso dQ\times_{TB}TE
 \qquad (r:E\to B).
 \label{eq:display}
\end{equation}
The given one-sided content unit and multiplication preserve the indicated
cartesian lifts, and the canonical cartesian--opcartesian mate is invertible
for pullbacks with display verticals and \emph{arbitrary} horizontals.
The one-sided contentad is already valid; no construction here turns an
arbitrary family of slice categories into one. The rig $\V$ has coherent
invertible distributors and absorption for $(\oplus,0),(\otimes,I)$; no
coproduct or product universal property is assumed. Recovery is on the
original categories and projections up to the specified canonical
reindexings, not a bijection of isomorphism classes. Display will prove
body equations only; grade and coefficient maps require separate arguments.

\section{The coefficient category and what bispans actually evaluate}

For $L\xleftarrow{s}U\xrightarrow{q}D\xrightarrow{t}K$ in finite sets, put
\begin{equation}
 F(L)=\V^L,\qquad
 (F(s,q,t)V)_k=
 \bigoplus_{d\in t^{-1}k}\ \bigotimes_{u\in q^{-1}d} V_{s(u)} .
 \label{eq:raw}
\end{equation}
This is a functor of \emph{categories}: it applies the indicated sums
and tensors to arbitrary tuples of $\V$-arrows. The isomorphism data of
the source bispans permute terms and their positions. To see the
pseudofunctoriality concretely, substitute the formula for one polynomial
into another. A term in the expansion is a target summand and a choice of
one inner summand for each of its outer positions; its positions are the
disjoint union of the chosen inner position sets. These are exactly the
summands and positions of the composite polynomial. Finite distribution
followed by tensor regrouping gives the compositor. On a triple
substitution the summands on both sides are the same addressed
three-level choices, and the positions are the same leaves in the
same specified order up to the same permutation. The compositor pentagon
is distributivity coherence for that tagged expansion. Formally, before
identifying equal input labels, assign separate variables to the primitive
occurrences in the unexpanded expression; nullary nonzero branches can be
tagged by fresh unit variables and then specialized to $I$. Nonzero
expanded monomials are distinct and squarefree in the formal variables;
compare the specified permutation of terms and factors, not expressions
with the same unlabelled endpoints. Empty sums and products use absorption
and unit coherence. Naturality moves nonstructural arrows through the
comparators. The singleton expression gives the unit triangle. This is
the finite-index version of the regular-coherence argument in
\cite[Sec.~3, structural coherence]{O1}.

Disjoint union of input sets goes to product of categories. Thus the
ordinary pseudofunctor on the $(2,1)$-category of finite bispans, with
target functors and natural isomorphisms between functors, gives after
taking nerves a product-preserving model in $\Cat_\infty$ (enlarge
universes for $\V$). Cnossen--Haugseng--Lenz--Linskens
\cite[Theorem A, Secs.~2.2, 3.3, 4.3]{CHLL} identify such models with
commutative semiring objects there. This is an explicit way to reach the
relevant side of their equivalence from the rig operations; it does not
assume that a distributive rig structure is a pair of adjunctions.
The formulas below are finite \emph{over each output point}, even if the
ambient bases are infinite: each output uses a finite bispan and
restriction to its finite list of input occurrences. Neither an infinite
coefficient sum nor a global assertion that $dQ$ is a finite set is used.

\paragraph{The augmentation is not a product-preserving change of model.}
For a carrier map $h:X\to Y$ write $h_!V$ for the additive transfer,
$(h_!V)_y=\bigoplus_{hx=y}V_x$, and $S_XV=\bigoplus_X V_x$.
Let $\Rcat=\int_{X\in\Fin} \V^X$ use the \emph{covariant} transfer
restriction of \eqref{eq:raw}. Its arrows are $(h,\alpha:h_!V\to W)$;
composition uses partition of inverse images. The comparison
$\rho_h:S_X V\iso S_Y h_!V$ defines a functor
$S:\Rcat\to\V$ sending $(h,\alpha)$ to $S_Y(\alpha)\rho_h$. There is an
identification over $\Fin$
\begin{equation}
 \Acat=\Rcat\times_{\V}\V_{/I}
      \ \simeq\ \int_{X\in\Fin}(S_X\downarrow I),
 \label{eq:global}
\end{equation}
where $\V_{/I}\to\V$ is domain. Objects have one arrow
$\epsilon:S_XV\to I$; arrows satisfy
$\epsilon'S_Y(\alpha)\rho_h=\epsilon$. This follows immediately on
objects and arrows, including compositions. No maps from separate
summands into $I$ have been extracted. Transfer to $Y$ carries the
augmentation by $\rho_h^{-1}$; it is the transfer subcategory on
which these total evaluations are pseudonatural.

In the discrete rig $(\mathbb N,+,\cdot)$, $(S_1\downarrow1)$ has one
object, whereas $(S_2\downarrow1)$ has two discrete objects, $(1,0)$
and $(0,1)$. Thus $X\mapsto(S_X\downarrow I)$ with these fibres
\emph{cannot} be a product-preserving bispan evaluator in general.
Nor does an arbitrary cartesian map of carrier polynomials yield a
coefficient arrow: the inclusion of one constant term into two
evaluates here to the nonexistent arrow $1\to2$. This is not the
claim that a $(2,1)$ source forces the \emph{values} to be groupoids.
Internal coefficient arrows in \eqref{eq:raw} and \eqref{eq:global}
can certainly be noninvertible.

\paragraph{Lifting an occurrence, including carrier-changing maps.}
For finite addressed $F$ and finite $X_i$, let
$K_X(V)(x)=\bigotimes_{i\in F} V_{i,x_i}$ on $\prod_i X_i$.
This is \eqref{eq:raw} with input $\coprod_i X_i$, one monomial
per output tuple, and its $F$ coordinates as positions.
For $h_i:X_i\to X'_i$, put $H=\prod_i h_i$. There is an isomorphism
of functors of the coefficient variables
\begin{equation}
 \Theta_h:H_!K_X(V)\iso K_{X'}((h_i)_!V_i),\quad
 (\Theta_h)_{x'}:
 \bigoplus_{x\in\prod_i h_i^{-1}(x'_i)}\bigotimes_i V_{i,x_i}
 \iso\bigotimes_i\bigoplus_{h_ix_i=x'_i}V_{i,x_i}.
 \label{eq:transfer}
\end{equation}
This is the inverse finite distributor, or the compositor for the two
factorizations of this indexed polynomial. It is not an arrow between a
polynomial and its image under an arbitrary map of shapes. Its compositional
equation, after partition comparators, is
\[
 \Theta_{kh}=(\Theta_k)_{h_!V}\, (\prod_i k_i)_!(\Theta_h)_V .
\]
Indeed both sides partition the same target fibre first by its intermediate
tuples and then by individual addresses. Its identity equation is the
singleton-fibre comparison. Thus a native family of arrows
$(h_i,\alpha_i)$ is sent at $x'$ to
$(\bigotimes_i\alpha_{i,x'_i})(\Theta_h)_{x'}$.
This is the necessary \emph{directed} lifting step.
The total $S K_X(V)\cong\bigotimes_i S_{X_i}V_i$ is natural for these
arrows by the same partition equation. Compose with $\bigotimes_i\epsilon_i$
and $\bigotimes_i I\cong I$. This defines the occurrence functor
$\bx_i:\Acat^F\to\Acat$ and not just its unaugmented part. For
$F=\varnothing$ the output is the singleton with $(I,\id_I)$.
An empty inverse image uses an actual map $0\to V'$, if present, and
absorption; no uniqueness of that map is presumed.

\paragraph{Dependent accumulation is not a simultaneous transfer.}
For $(X,V,\epsilon)$ and continuations $(Z_x,W_{x,-},\delta_x)$,
the coefficient at $(x,z)$ of their composite $P\St D$ is
$V_x\otimes W_{x,z}$. Total augmentation first contracts each $W_{x,-}$
by $\delta_x$ and then the remaining $V$ by $\epsilon$.
For a native map $(h,\alpha)$ and $(g_x,\beta_x):D_x\to D'_{hx}$,
the fibre over $(x',z')$ has the following map:
\begin{align}
 \bigoplus_{hx=x'}\bigoplus_{g_x z=z'} V_x\otimes W_{x,z}
 &\cong\bigoplus_{hx=x'} V_x\otimes
             \Big(\bigoplus_{g_xz=z'} W_{x,z}\Big)
 \xrightarrow{\oplus_x(1\otimes\beta_{x,z'})}
       \bigoplus_{hx=x'}V_x\otimes W'_{x',z'}
 \notag\\[-2pt]
 &\cong\Big(\bigoplus_{hx=x'}V_x\Big)\otimes W'_{x',z'}
 \xrightarrow{\alpha_{x'}\otimes1}V'_{x'}\otimes W'_{x',z'} .
 \label{eq:dep}
\end{align}
The first stage has a separate inner transfer at each $x$. In the second,
$W'_{x',z'}$ is pulled back from the \emph{target}, hence constant on the
fibre of $h$. Taking separate sums of all outer and all inner variables
before tensoring would introduce terms from \emph{different} ancestors.
For a subsequent continuation map, the inner operation at $x'$ is
constant on the fibre already merged at $x'$, so it commutes with the
previous outer aggregation by tensor bifunctoriality, after distributing
the intermediate sums. This proves compatibility with composition; the
two incoming augmentation equations, used inner first, prove the total
augmentation condition. For zero fibres the same identity is naturality
of the absorption isomorphism, not cancellation out of an initial object.

For the comparison
\begin{equation}
 \bx_i(P_i\St D_i)\iso (\bx_i P_i)\St
               \bigl(x\longmapsto\bx_i D_{i,x_i}\bigr) ,
 \label{eq:Fubini}
\end{equation}
the carrier sends $((x_i,z_i))$ to $((x_i),(z_i))$ and the coefficient
is the permutation collecting the $V$ factors ahead of the $W$ factors.
At a fixed target tuple, both native-map composites expand to the tensor
over $i$ of \eqref{eq:dep}, conjugated by that shuffle. In particular
each $\beta_{i,x_i,z'_i}$ precedes $\alpha_{i,x'_i}$. All extra exchanges
are distributor naturality or act on disjoint tensor factors. The
augmentations have the same inner-before-outer order. This proves
naturality in the actual aggregation maps. The same tagged expansion proves
compatibility with reassociation, permutation and the two units.
This concrete interface lemma (cf.\ \cite[full native dependent
Fubini]{O1}) uses no extra algebraic axiom on $\V$. It \emph{does}
use the Grothendieck arrow definition and its order of evaluation; that
definition is not the conclusion of the bispan classification theorem.

\section{The relative corners calculation in the content double category}

Let $\mathbf N$ denote the double category of contentful arrows.
Its horizontal arrow $n=(Q,a):A\rightsquigarrow B$ consists of
$Q\in\Lcat_B$ and $a:A\to dQ$. A square to $(Q',a')$ with verticals
$k:A\to A'$ and $u:B\to B'$ consists of a total grade arrow
$m:Q\to Q'$ over $u$ with $dm\,a=a'k$.
This is equivalently a grade arrow $u_!Q\to Q'$ in the target fibre.
The horizontal composite of $(Q,a)$ followed by $(R,b):B\rightsquigarrow C$
has grade $R\diamond b_!Q$ and body $\mu^L_{R,b_!Q}\,d\kappa_{b,Q}\,a$.

\begin{lemma}[A relative corners square, stable under actual composition]
\label{lem:corner}
Let $r:E\to B$ lie in $\J$. For $(Q,a)$ as above, take
\[
 A_r=A\times_{dQ}d(r^*Q),\quad
 \ell:A_r\to A,\quad
 a_r:A_r\to d(r^*Q).
\]
The square with top $(r^*Q,a_r)$, bottom $(Q,a)$ and verticals $(\ell,r)$
is cartesian for $\mathrm{cod}:\mathbf N_1\to\Set$ against all squares.
One has $\ell\in\J$. Lifting through displays is stable, up to the
canonical cartesian comparisons, for identities and both vertical and
horizontal pasting. No lifts outside $\J$ are used.
\end{lemma}
\begin{proof}
Suppose another square into the bottom arrow has grade $m:Z\to Q$
over $rv$ and left edge $l:A_0\to A$, with source body $z:A_0\to dZ$.
The cartesian property of $c_{r,Q}$ gives a unique $m':Z\to r^*Q$
over $v$ with $c_{r,Q}m'=m$. The pullback gives the unique left
edge $\langle l,dm'\,z\rangle$ of the factor square. This is exactly
codomain-cartesianness; $v$ is unrestricted. Display identifies
$A_r$ with $A\times_{TB}TE$ via $\nu_Q a$. Thus $\ell$ is a
pullback of $Tr\in\J$. Vertical composition is composition of
cartesian lifts. For horizontal identity the grade arrow is $U(r)$,
cartesian by hypothesis, and the source pullback is $E$ over $B$
by cartesianness of $\eta^T$ and $\nu_{U_B}\eta^L_B=\eta^T_B$.

For horizontal composition take $n=(Q,a)$ followed by $m=(R,b)$ as
above, and a display $r:C'\to C$. Put $R'=r^*R$,
$j=d c_{r,R}:dR'\to dR$ and
$B'=B\times_{dR}dR'$ with projections $v:B'\to B$, $b':B'\to dR'$.
Both $j$ and $v$ belong to $\J$. The grade of the pasted top arrow is
canonically identified as
\[
 R'\diamond b'_!v^*Q
 \ \cong\ R'\diamond j^*b_!Q
 \ \cong\ r^*(R\diamond b_!Q).
\]
Here the first is the BC mate for $(v,j;b,b')$ and the second the
cartesian multiplication comparison. They identify the pasted grade
arrow \emph{over the original grade} with its cartesian lift, by their
factorization equations. For the left source, display gives
$B'\cong B\times_{TC}TC'$ using $\nu_R b$. The twice-lifted source is
\begin{equation}
 A\times_{TB}TB'\ \cong\ A\times_{TTC}TTC'
 \ \cong\ A\times_{TC}TC',
 \qquad \delta=T(\nu_R b)\,\nu_Q a:A\to TTC .
 \label{eq:composegeom}
\end{equation}
The first comparison uses preservation of pullbacks by $T$, the
second the cartesian naturality square for $\mu^T$ at $r$. The actual
composite body has occurrence $\mu^T_C\delta$: this uses naturality
of $\nu$ at $\kappa_{b,Q}$ and its multiplication equation. In these
coordinates the pasted comparison sends a two-level decoration to
its $\mu^T$-flattening, by naturality of the actual bodies and their
arity equation. Hence \eqref{eq:composegeom} identifies the pasted
square with the direct cartesian square, not just their source sets.
Neither preservation of opcartesian arrows by $d$ nor invertibility
of $\mu^L$ is required.
\end{proof}

This lemma has a literal double-category interpretation. In the
horizontal opposite of $\Sq(\Set)$ restrict \emph{horizontal} arrows
to backwards $\J$-maps and keep arbitrary vertical maps and all
commuting squares. A content arrow followed by a backward display is
sent to a backward display followed by the content arrow from the lemma.
For a map of such pairs, paste its input content square with the
first lift; the second lift is cartesian, so the result factors uniquely
over the upper map of the display square. This defines the output square.
Uniqueness gives identity and composition for maps of pairs.
The unit and multiplication comparisons for successive backwards maps
are vertical composition of cartesian squares; those for successive
content arrows are the horizontal identity/composition parts of the lemma.
Their pasting equations hold in this corner calculus: on pasting with
the common chosen cartesian comparison to the original arrow, both sides
are the same pasting of input squares, using horizontal associativity and
units of $\mathbf N$. The universal property makes the two factorizations
equal. This is the proof mechanism of \cite[Def.~6.1, Prop.~6.3]{BG},
restricted where lifts are taken, not where they are tested. Under
horizontal reversal its variance is the usual corners variance.
The ambient product of formally composable arrows in
\cite[Def.~3.2 and Sec.~3.2]{BG} is not yet the relative
pseudocoend, and a backward display alone is not the context module.

\paragraph{Why a vertical restriction would be wrong.}
Already in the square double category of sets, keep only surjective
vertical arrows and commuting squares with both verticals surjective.
Consider the horizontal $u:2\to1$ and the prescribed right edge
$f:2\to1$, with $2=\{0,1\}$. For each constant $s_i:2\to2$ there is
a square from $s_i$ to $u$ with verticals $(\id_2,f)$. If a
cod-cartesian lift had horizontal $h:R\to 2$ and left projection
$l:R\twoheadrightarrow 2$, each square would factor with left
$r_i:2\twoheadrightarrow R$, $l r_i=\id$, $h r_i=s_i$. Surjectivity
of $r_i$ forces $h$ to be constantly $i$ for both $i$, a contradiction.
A class closed under pullback/composition is therefore no licence to
apply the global result in this restricted vertical category.

\section{Assembly and comparison on the actual relative tensor}

Let $P\in\Kcat_B$ have augmented fibres $(E_b,V_{b,-},\epsilon_b)$,
total carrier $E$ and display $\pi:E\to B$ in $\J$.
A contextful arrow is the decorated span $B\xleftarrow{\pi}E\xrightarrow{g}C$,
not just its backwards first leg. Context composition forms dependent
pairs with the coefficients of \eqref{eq:dep}. For each $Q\in\Lcat_B$
put $D=dQ\times_{TB}TE$. At $y\in dQ$ with $\nu_Q(y)=(s,(b_i))$,
$D_y=\prod_{i\in E_s}E_{b_i}$. Repeated labels at distinct positions
are independent factors. Their single-cut formula is
\begin{equation}
 \widehat P_y=\bx_{i\in E_s}P_{b_i},\qquad
 \widehat Q=\pi^*Q,\qquad
 \chi_{P,Q}:e\widehat P\cong D\xrightarrow{D_{\pi,Q}^{-1}}d(\pi^*Q).
 \label{eq:cut}
\end{equation}
Here $e\widehat P\to dQ$ is a pullback of $T\pi$ and is in $\J$.
More explicitly, the raw functor sending $(V_e)_{e\in E}$ to the
coefficients on $D$ is the pointwise evaluator for
\[
 E\ \longleftarrow\ \{(z,i):z=(y,t)\in D,\ i\in E_{s(y)}\}
 \ \longrightarrow\ D\ \xrightarrow{\id}\ D,
 \quad (z,i)\longmapsto t(i)\ \text{on the left}.
\]
At each $y$, apply the total augmentation after the independent
transfer calculation. The chosen display of the \emph{same} $D$ supplies
the last arrow of \eqref{eq:cut}. For general representatives
$(Q,a:A\to dQ)$ and $(P,g:E\to C)$ the result is
\[
 \big(a^*\widehat P,\quad g_!\pi^*Q,\quad
 e(a^*\widehat P)\to e\widehat P\xrightarrow{\chi_{P,Q}}d(\pi^*Q)
                    \xrightarrow{d\kappa_{g,\pi^*Q}}d(g_!\pi^*Q)\big).
\]
Its carrier source is the $A_\pi$ of Lemma~\ref{lem:corner}.

\paragraph{A directed square is not a reindexing isomorphism.}
Consider $\pi':E'\to B'$, $\pi'h=u\pi$, $m:Q\to Q'$ over $u$,
and a native context map over $u$ with total carrier $h:E\to E'$.
Write $\widetilde m:\pi^*Q\to\pi'^*Q'$ for the unique arrow over $h$
with $c_{\pi',Q'}\widetilde m=m c_{\pi,Q}$. Its factorization is
\begin{equation}
 \widetilde m=\lambda_{h,m}\kappa_{h,\pi^*Q},\qquad
 \lambda_{h,m}:h_!\pi^*Q\longrightarrow\pi'^*Q'.
 \label{eq:grade}
\end{equation}
Both constructions are legitimate for $h\notin\J$; neither invokes $h^*$.
This arrow need not be invertible. At the other end, $\nu_{Q'}dm=Tu\,\nu_Q$
matches the shape and carries each label $b_i$ to $ub_i$. On its tuple
fibre use \eqref{eq:transfer} followed by the tensor of the given
aggregation maps. This includes an arbitrary total grade map $m$,
not just an invertible display square. The comma body equation is
\[
 d\widetilde m\,\chi_{P,Q}
    =\chi_{P',Q'}(dm,Th):D\longrightarrow d(\pi'^*Q').
\]
The first coordinate follows from the defining grade factorization and
the second from naturality of $\nu$ on $\widetilde m$. Conclude from
\eqref{eq:display} only after checking the native maps. For two squares,
$\widetilde {m'm}=\widetilde m'\widetilde m$ by cartesian uniqueness;
opcartesian factorization gives the corresponding composition equation
for $\lambda$. Equation \eqref{eq:transfer} gives the coefficient
composition equation. Reindexing at the source and transporting the
output grade give the moving endpoint actions.

\paragraph{Why every pure middle arrow matters.}
Take any $u:B\to B'$, any admitted $\pi':E'\to B'$ and its pullback
$(\pi_u:E_u\to B, v:E_u\to E')$. Passing $u$ across the cut changes
$Q$ to $u_!Q$ on one side and $P'$ to $u^*P'$ on the other. On the
source representative the fibres and coefficients identify by label
substitution. The necessary grade comparison is exactly
\begin{equation}
 \lambda_{v,\kappa_{u,Q}}:
 v_!\pi_u^*Q\iso\pi'^*u_!Q .
 \label{eq:balance}
\end{equation}
Indeed \eqref{eq:grade} is the defining mate equation. Before
postcomposition by a later $d\kappa_g$, the two body routes have
coordinates $(d\kappa_{u,Q}(y),Tv(t))$ in
$d(u_!Q)\times_{TB'}TE'$. Thus neither route assumes $d\kappa_u$
cartesian. For successive pure maps the two applications of
\eqref{eq:balance}, after opcleavage comparison, have the same
composite against the original cartesian lift as the single mate.
Cartesian and opcartesian uniqueness give $\BC_{u'u}=\BC_{u'}\BC_u$
with the indicated whiskerings, and $\BC_{\id}=1$. The same
factorization proves naturality for \eqref{eq:grade}, and the
coefficient squares are the pointwise naturality of \eqref{eq:transfer}.
This checks the generators and relations of the pseudobalanced
coend, not just the commuting squares that defined the relative
corners operation.

\begin{theorem}[Concrete recovery, not an abstract comparison theorem]
\label{th:assembly}
With the above choices, the decorated relative-corners calculation gives
the same pseudonatural crossing on $N\odot M$ as the independent
interaction law of \cite{O1}, after the standard normalization of
$M\odot N$. Its four canonical structure maps and the corresponding
composite on $H(A,B)=(e_A\downarrow d_B)$ agree with the original
ones. No such assertion holds simply by evaluating a functor on the
unadorned square or polynomial diagrams.
\end{theorem}
\begin{proof}
A pseudocoend here is presented by native product arrows and invertible
balance generators modulo their identity, composition, and naturality
equations. Equations \eqref{eq:cut}--\eqref{eq:balance} give the same
objects and images of both kinds of generators as the independent law.
For $M\odot N\simeq H$, normalization sends
$[Z;(P,b:eP\to Z),(Q,c:Z\to dQ)]$ to $(P,Q,cb)$. The inverse
uses $Z=eP$. On a comma arrow $(\alpha,\beta)$ the inverse applies
$\beta$, balances along $e\alpha$, then applies $\alpha$.
This proves retention of the native arrows, including two arrows with
equal $e\alpha$; the balance along $b$ furnishes the comparison to
the original representative \cite[section on pseudocoend normalization]{U}.

For clarity the four constraints are as follows, with canonical
associators implicitly inserted:
\begin{alignat*}{2}
 \tau(j_NM)&\cong Mj_N:
 &\quad &\text{one position on }\eta^L_B(B),\quad U_E\cong\pi^*U_B;\\
 \tau(Nj_M)&\cong j_MN:
 &&\text{forced unit choices},\quad (\id_B)^*Q\cong Q;\\
 \tau(m_N M)&\cong (M m_N)(\tau N)(N\tau):
 &&\text{flatten addressed leaves},\quad
    \pi^*Q\diamond (dc_{\pi,Q})^*R\cong\pi^*(Q\diamond R);\\
 \tau(N m_M)&\cong (m_M N)(M\tau)(\tau M):
 &&\text{dependent Fubini},\quad
             (\pi\rho)^*Q\cong\rho^*\pi^*Q .
\end{alignat*}
Here $\rho$ is the second context projection and $R\in\Lcat_{dQ}$.
The grade comparisons are the cartesian factors over the original
grades or composite cleavage. The coefficient comparisons are the
singleton unitor, unit-product comparison, addressed regrouping, and
\eqref{eq:Fubini}, respectively. At the first and third constraints the
source is restricted along the actual $\eta^L_B$ and $\mu^L_{Q,R}$;
the first display coordinate in the multiplication case is
$\mu^L_{Q,R}(z)$, not $z$. The decorated second coordinate is the
flattened $T$-label; cartesianness of $\mu^T$ retains the leaf bijection.
Display proves the structural body squares. These comparisons agree as
\emph{actual} components with those of \cite[construction of the full
interaction]{O1}, so all ten equations transfer. This last inference is
by identification on generators and modifications, \emph{not} by
identification on carriers or by applying a new general coherence theorem.

For a second check, the following ledger gives the common coefficient
calculation on each source pasting; the grade/body checks for all rows
are as above. It uses the numbering of \cite[Def.~4.1, App.~A]{GL}.
\begin{center}\small
\begin{tabular}{@{}llp{.55\linewidth}@{}}
\toprule
& Comparison & Witness at the source representative\\
\midrule
C1&three content stages&Same flattened leaves $(k,j,i)$ and tensor regrouping.\\
C2&content left unit&Remove the left singleton-position layer.\\
C3&content multiplication/context unit&All choices forced, all primitive factors $I$.\\
C4&two units&One unit position with a forced choice.\\
C5&two multiplications&At each addressed leaf $(j,i)$, split $(x_{ji},z_{ji})$; splitting commutes with flattening by \eqref{eq:Fubini}.\\
C6&context multiplication/content unit&At the one unit position retain the dependent pair.\\
C7&three context stages&Same ancestor chain $(x,z,w)$ and tensor factors.\\
C8&context right unit&Remove the last forced choice and its $I$.\\
C9&content right unit&Remove the trailing singleton-position layer.\\
C10&context left unit&Remove the initial forced choice and its $I$.\\
\bottomrule
\end{tabular}\end{center}
In C5 both native routes perform the $x_{ji}$-dependent continuation
maps before the corresponding outer merges. The rest are the
constant-on-the-merged-fibre exchanges of \eqref{eq:dep}. In C1/C2/C9
the polynomial monad laws identify position maps, not merely arity
cardinalities. On grades, compose the two routes with the same
cartesian arrow to the original grade: the one-sided law gives the
same total arrow, with opcartesian factorization and mate pasting at
moving endpoints. Genuine cartesian uniqueness identifies the factors.
Only then do the two coordinates identify the body. Thus ambient
associators of balanced representatives, native arrows and interchangers
have the same images in the normalized homs.

It follows that $(j_M\odot j_N)$ and
$(m_M\odot m_N)(M\odot\tau\odot N)$ yield the original relative
pseudopromonad after normalization, by the composition theorem used
in \cite[section on mixed algebra]{U}. For $(P,Q,f)$ then $(D,S,g)$
the output grades are $P\St f^*\widehat D$ and
$S\diamond g_!\widehat Q$. In particular the latter is the whole grade,
not its restriction to the image of the evaluated body.
\end{proof}

\section{Separating tests and exact boundary}

\begin{example}[A lawful relative calculation with no global co-crossing]
Let $G(B)=\varnothing$ for $B=\varnothing$, and $G(B)=1$ otherwise.
Use the constant discrete grades $\id,G$ with substitution by functor
composition, extension by evaluation and $T=\mathsf{List}$.
The occurrence of the identity grade is the singleton list; that of $G$
is the empty list wherever defined. The equations for substitution follow
from singleton/empty-list concatenation and $G^2=G$. Cartesianness of
the grade algebra and BC are automatic for constant discrete transport.
For a surjection $r:E\to B$, the two sets have the same inhabitation
status, and the only word above the empty word is empty. Thus
\eqref{eq:display} holds, including $\varnothing\to\varnothing$.
For $\varnothing\to1$ the display is $\varnothing\to1$ and fails.

The corresponding global-co-crossing failure is a horizontal-pasting
failure even though the grade fibration here has global lifts.
Let $i_0,i_1:1\to2$ have disjoint images. Take the $G$ arrow
$1\rightsquigarrow1$ with unique body, followed by the pure $i_0$.
Their composite has grade $G$ and body $1\to G(2)$. Pulling the pure
$i_0$ back along $i_1$ first gives source $\varnothing$; pulling the
$G$ arrow back at this intermediate $\varnothing\to1$ still gives
source $\varnothing$. Pulling the composite back along $i_1$ directly
instead gives $1\times_{G(2)}G(1)=1$. The two source sets are not
isomorphic. Thus global co-crossedness in \cite[Def.~6.1]{BG} cannot be
silently imposed, even though the calculus over finite-fibre surjections
is lawful, with \emph{all} maps still present as pure endpoint changes.
\end{example}

\begin{example}[The same carrier can carry two coefficient maps and a grade fold]
Over $k=\mathbb Q$, set $V=k^2$, $\epsilon(a,b)=a+b$, and
$p(a,b)=(a+b,0)$. A source with two $k$-rows and total sum augmentation
maps over the carrier fold $h:2\to1$ to the singleton with $(V,\epsilon)$.
The two aggregation components $\id_V$ and $p$ both preserve augmentation.
At two occurrences of the same base label, \eqref{eq:transfer} sends them,
after identifying four source rows with $V\otimes V$, to
$\id_{V\otimes V}$ and $p\otimes p$. On
$e_2\otimes e_2$ the latter gives $e_1\otimes e_1$, not $e_2\otimes e_2$.
The same difference occurs for the two singleton endomorphisms.
Moreover the singleton categorical product over $k$ has dimension
$\dim(V\times_k V)=3$, whereas $\dim(V\otimes V)=4$.

The content part can be noninvertible in this very same fold.
Use the valid indexed-receipt content of \cite[Indexed receipts]{O1}: a grade at $B$
is a set of tags $q:Q_0\to TB$, $dQ=B+Q_0$, reindexing pulls tags
back along $T\pi$, and pushforward postcomposes the labels. At $B=1$
choose one tag labelled $[*,*]$. For $\pi:2\to1$ and
$\pi':1\to1$, $\pi^*Q$ has four tags, one per decorated word.
$h_!\pi^*Q$ still has four tags, now all at $[*,*]$;
$\pi'^*Q$ has only one. Thus \eqref{eq:grade} is the genuine
grade fold $4\to1$, not an isomorphism of labels. Both choices of
coefficient map have precisely this same grade fold. The comparison
must distinguish the coefficients \emph{and} retain the directed grade
fold. Neither the unlabelled square factor nor a passage to the cores
of the target categories meets that requirement.
\end{example}

\begin{example}[Surjective linear maps, without an additive adjunction]
Take finite-dimensional $k$-spaces and only surjective linear maps,
with tensor and direct sum. Over a field both operations preserve
surjections; the rig comparators are isomorphisms. There is no
$0\to k$, so no empty carrier is augmentable. Each nonempty finite
carrier is augmentable (one $k$, remaining rows zero); hence every
admitted finite-fibre surjection is realizable in ordinary sets with
choice. The relative $G$ example therefore applies in a nonthin rig
with noninvertible native maps, e.g.\ the augmentation
$(k^2,\epsilon)\to(k,\id_k)$ over a singleton. The previous
$k^2$-endomorphism $p$ is \emph{not} claimed surjective and is not used
as a witness in this subcategory.

Direct sum cannot be the categorical coproduct of $k,k$ in this
category: there is no surjection $k\to k\oplus k$. Stronger, the
diagonal need not admit any left adjoint: for $k=\mathbb Q$, a
coproduct cone on $k,k$ with surjective legs would have target of
dimension at most one and must have dimension one to receive a
pair of maps to $k$ by factorization. Fixing the first nonzero scalar
then fixes the mediator, although the second scalar can vary. More
generally the pair $k^2,k^2$ cannot have a coproduct over \emph{any}
field: the pair of identities into $k^2$ forces surjections from and
to $k^2$ for the proposed coproduct object, hence both legs are
isomorphisms. Fixing the first component to be identity forces the
mediator, whereas the second component can be varied through at least
two elements of $\mathrm{GL}_2(k)$. Thus the fold transfer here is a
perfectly usable rig sum, not the left adjoint to reindexing. This
distinction also prevents an accidental assertion about the number of
nonzero scalars over $\mathbb F_2$.
\end{example}

\paragraph{Recognition uses more than construction.}
Let $\mathcal D$ be the realized context projections. For a marked
crossing, the claimed exhaustive carrier $D$, invertible $\chi$ and
preservation of the two coordinates imply $D_{\pi,Q}\chi=1$ and
hence $D_{\pi,Q}=\chi^{-1}$. Without invertibility this yields only a
section. Preservation of the original unit and multiplication grade
projections forces their components to factor the given arrows
$U(\pi)$ and the accumulation of $(c_{\pi,Q},c_{j,R})$ through the
respective cartesian lifts; uniqueness forces the canonical factors.
Invertibility gives cartesian algebra. The marked pure-middle square
likewise forces the actual mate \eqref{eq:balance}. This is the
recognition argument of \cite[marked recognition]{O1}; it tests only
$\mathcal D$ and its actual pullbacks. For the augmented rig considered here, under choice $\mathcal D$ is
the finite-fibre part of $\J$ if $\V(0,I)$ is inhabited, otherwise its
finite-fibre surjections. Every nonempty finite fibre is realized by
putting $I$ on one chosen element and $0$ on the rest, while an empty
fibre is realizable exactly when $0\to I$ exists.
The CHLL equivalence neither names these marks nor quantifies over
interactions; BG's recognition of calibrations \cite[Thm.~8.2]{BG}
is a different statement.

\paragraph{What remains outside this comparison.}
A theorem identifying an arbitrary relative double-category corners
construction, after a functorial decoration, with a law of relative
modules would require an ambient comparison sending the product of
formally composable cell categories to the balanced module tensor,
extending backward displays by the unrestricted forward tails and
carrying the slice construction, its directed transfers and dependent
multiplication. Equations \eqref{eq:transfer}, \eqref{eq:dep},
\eqref{eq:Fubini}, \eqref{eq:grade} and \eqref{eq:balance}, together
with the four structural identifications, are concrete tests for such
a functor. They have been checked here for this presentation. No
universal, unmarked, or model-independent such functor has been
constructed. BG's ordinary polynomial identification
\cite[Props.~10.2--10.3]{BG} uses a calibration into spans;
it is not that functor. Capucci--Myers \cite[Def.~D.9]{CM} give a
display-span wreath over a pseudocategory with an intertwiner as
part of its data, not an independent derivation of this intertwiner
on the augmented native categories.

\begin{thebibliography}{9}
\bibitem{O1} \href{https://github.com/eternal-november/eos-research-programmes/blob/f458186590df968a516ba73a8eb76fcf044544b0/projects/contextual-composition/interaction/independent-criterion/General_Independent_Interaction_Criterion.tex}{\emph{General independent interaction criterion}}.
Project source. Sections on native dependent Fubini, full interaction,
and marked recognition.
\bibitem{U} \href{https://github.com/eternal-november/eos-research-programmes/blob/f458186590df968a516ba73a8eb76fcf044544b0/projects/contextual-composition/foundations/relative-construction/delivery/One_relative_construction_for_context_content_and_effects.tex}{\emph{One relative construction for context, content and effects}}.
Project source. Sections on pseudocoend normalization and mixed algebra.
\bibitem{O8} \href{https://github.com/eternal-november/eos-research-programmes/blob/f458186590df968a516ba73a8eb76fcf044544b0/projects/contextual-composition/synthesis/Integrated_Relative_Computation_Acceptance_Dossier.tex}{\emph{Integrated relative computation acceptance dossier}}.
Project source. Section on independent inputs and the coherent
interaction; crosswalk to the detailed O1 proof. All three links resolve to
the same pinned repository revision, \texttt{f458186590df}.
\bibitem{BG} M. Buckley and R. Garner,
\emph{Composing Double Categories},
\href{https://arxiv.org/html/2610.41551v2}{arXiv:2610.41551v2},
Secs.~3, 6, 8, 10.
\bibitem{CHLL} B. Cnossen, R. Haugseng, T. Lenz, and S. Linskens,
\emph{Homotopical commutative rings and bispans},
\href{https://arxiv.org/html/2403.06911v2}{arXiv:2403.06911v2},
Theorem A; Def.~2.2.4, Cor.~3.3.7, Prop.~4.3.1.
\bibitem{CM} M. Capucci and D. J. Myers,
\emph{Contextful Arrows via Wreath Products},
\href{https://arxiv.org/pdf/2410.21889v2}{arXiv:2410.21889v2},
App.~D, Def.~D.9.
\bibitem{GL} N. Gambino and G. Lobbia,
\emph{On the formal theory of pseudomonads and pseudodistributive laws},
\href{https://arxiv.org/pdf/0907.1359v3}{arXiv:0907.1359v3},
Def.~4.1, Remark~4.3 and App.~A.
\end{thebibliography}
\end{document}
